\SetPractica{Humedad y cenizas}{Determinación de humedad y cenizas}

\section{Introducción teórica}



\section{Objetivos}

\begin{itemize}
    \item 
\end{itemize}

\section{Materiales, equipos y reactivos}

\begin{table}[H]
\centering{\small\setstretch{1.25}
\begin{tabular}{|m{0.45\linewidth}|m{0.45\linewidth}|} \hline
    
    \thead{Equipos} & \thead{Materiales} \\ \hline
    
    {Horno de secado \par}
    {Balanza analítica \par}
    {Campana de extracción de gases \par}
    {Mufla \par}
    {Placa de calentamiento}
    
    &
    
    {Cápsulas de porcelana \par}
    {Desecador \par}
    {Pinzas \par}
    {Guantes para calor \par}
    {Espátula \par}
    {Mortero y pistilo}
    
    \\ \hline
\end{tabular}}\end{table}

\section{Procedimientos}

\subsection{Preparación de cápsulas}
\begin{enumerate}
    \item Introducir las cápsulas o crisol al horno a una temperatura de 105°C $\pm$ 2°C, 20min. como mínimo. 
    \item Trasladar la cápsula o crisol al desecador
    \item Dejar enfriar por 20min. como mínimo.
    \item Pesar las cápsulas o crisoles en balanza analítica y repetir el ciclo horno-desecador y hasta obtener una diferencia $\leq$ 0,0005g en dos pesadas consecutivas. 
    \item Registrar como M\textsubscript{1} considerando para los cálculos el último valor de la masa.
\end{enumerate}

\subsection{Medición de porcentaje de humedad}
\begin{enumerate}
    \item Se tritura en el mortero la muestra si es sólida.
    \item Se pesan de 3 a 5g de muestra en la cápsula que previamente fue puesta a peso constante (M\textsubscript{2}).
    \item Evaporar a sequedad en el horno de secado a 105 °C $\pm$ 2 °C por 1h.
    \item Trasladar la cápsula al desecador y dejar enfriar por 20min. como mínimo. Llevar la cápsula a masa constante repitiendo el ciclo horno-desecador, hasta obtener una diferencia $\leq$ 0.0005g en dos pesadas consecutivas.
    \item Registrar como M\textsubscript{3}, la última masa obtenida.
\end{enumerate}

\subsection{Medición de porcentaje de cenizas}
\begin{enumerate}
    \item Registrar cómo M\textsubscript{2}, el valor de la cápsula a peso constante con la muestra tomada de la determinación de porcentaje de humedad.
    \item Colocar las cápsulas con muestra en placa de calentamiento para llevar a combustión total, hasta que deje de emitir humos (en campana de extracción).
    \item Dejar enfriar y llevar a la mufla a 550°C por 3h.
    \item Una vez fría la mufla, sacar las cápsulas y trasladar la cápsula al desecador y dejar por 20min. como mínimo. Llevar la cápsula a masa constante repitiendo el ciclo horno-desecador, hasta obtener una diferencia $\leq$ 0.0005g en dos pesadas consecutivas.
    \item Registrar como M\textsubscript{4}, la última masa obtenida.
\end{enumerate}

\subsection{Cálculo de porcentajes}

\subsubsection{Cálculo de porcentaje de humedad}
\begin{equation*}
    \text{Porcentaje de humedad} = 100 - \frac{M_{3} - M_{1} \cdot 100}{M_{2}}
\end{equation*}

Donde: 
\begin{itemize}
    \item M\textsubscript{3} es la masa de la cápsula con el residuo, después de la evaporación, en g;
    \item M\textsubscript{1} es la masa de la cápsula vacía a masa constante, en g, y
    \item M\textsubscript{2} es la masa de muestra, en g.
    \item 100 es el factor de porcentaje
\end{itemize}


\subsubsection{Cálculo de porcentaje de cenizas}
\begin{equation*}
    \text{Porcentaje de cenizas} = 100 - \frac{M_{4} - M_{1} \cdot 100}{M_{2}}
\end{equation*}

Donde: 
\begin{itemize}
    \item M\textsubscript{4} es la masa de la cápsula con el residuo, después de la calcinación, en g;
    \item M\textsubscript{1} es la masa de la cápsula vacía a masa constante, en g, y
    \item M\textsubscript{2} es la masa de muestra, en g.
    \item 100 es el factor de porcentaje
\end{itemize}